\section{编写 CUDA 内核}

\subsection{CUDA 编程基础}

CUDA 是 NVIDIA 的 C/C++ API，用于编程 GPU。核心思想很简单：写一个函数 $f$，CUDA 对所有元素并行执行 $f$。

\begin{figure}[H]
\centering
\begin{tikzpicture}[
    block/.style={draw, minimum width=1.8cm, minimum height=0.6cm, font=\small},
]
\node[block, fill=blue!10] (b00) at (0,0) {Block (0,0)};
\node[block, fill=blue!10] (b01) at (2.2,0) {Block (0,1)};
\node[block, fill=blue!10] (b02) at (4.4,0) {Block (0,2)};
\node[block, fill=green!10] (b10) at (0,-0.9) {Block (1,0)};
\node[block, fill=green!10] (b11) at (2.2,-0.9) {Block (1,1)};
\node[block, fill=green!10] (b12) at (4.4,-0.9) {Block (1,2)};
\node[font=\small, above=0.2cm of b01] {Grid: numBlocks = (2, 3)};
\node[font=\scriptsize, below=0.2cm of b11, text width=5cm, align=center] {每个 Block 包含 blockDim 个线程\\每个线程由 (blockIdx, threadIdx) 唯一标识};
\end{tikzpicture}
\caption{CUDA Grid-Block-Thread 组织结构}
\end{figure}

关键概念：
\begin{itemize}
    \item \texttt{blockIdx}：当前 Block 在 Grid 中的索引
    \item \texttt{blockDim}：每个 Block 的线程数
    \item \texttt{threadIdx}：当前线程在 Block 中的索引
    \item 全局线程索引：$i = \text{blockIdx.x} \times \text{blockDim.x} + \text{threadIdx.x}$
\end{itemize}

\subsection{GeLU 的 CUDA 实现}

\begin{lstlisting}[caption={GeLU CUDA 内核（gelu.cu）},language=C]
__global__ void gelu_kernel(
    const float* in, float* out, int num_elements
) {
    // 计算全局线程索引
    int i = blockIdx.x * blockDim.x + threadIdx.x;

    // 边界检查
    if (i < num_elements) {
        float x = in[i];
        float a = 0.79788456f * (x + 0.044715f * x * x * x);
        float tanh_a = tanhf(a);
        out[i] = 0.5f * x * (1.0f + tanh_a);
    }
}
\end{lstlisting}

\begin{knowledgebox}{CUDA 内核的关键模式}
几乎所有 CUDA 内核都遵循以下模式：
\begin{enumerate}
    \item 计算全局索引 $i$：\texttt{blockIdx.x * blockDim.x + threadIdx.x}
    \item 边界检查：\texttt{if (i < num\_elements)}——因为最后一个 Block 可能有多余的线程
    \item 执行计算：从输入指针读取、计算、写入输出指针
\end{enumerate}
\texttt{\_\_global\_\_} 关键字标识该函数为 CUDA 内核函数。
\end{knowledgebox}

\subsection{Python 端启动 CUDA 内核}

PyTorch 提供了 \texttt{load\_inline} 工具，可以在 Python 中直接编译和调用 CUDA 代码：

\begin{lstlisting}[caption={在 Python 中编译和使用 CUDA 代码}]
from torch.utils.cpp_extension import load_inline

# CUDA 源码（字符串形式）
cuda_src = open("gelu.cu").read()
cpp_src = "torch::Tensor gelu(torch::Tensor x);"

# 编译为 Python 模块
module = load_inline(
    cuda_sources=[cuda_src],
    cpp_sources=[cpp_src],
    functions=["gelu"],
    extra_cflags=["-O2"],
    name="inline_gelu",
)

cuda_gelu = module.gelu  # 现在可以像普通函数一样调用
\end{lstlisting}

\subsection{CUDA GeLU 性能}

Profiling CUDA GeLU 的结果显示，我们的自写内核确实只启动了\textbf{一个 CUDA 内核}（名为 \texttt{gelu\_kernel}），所有计算在内核内部完成。与 PyTorch 的差距来自于：

\begin{itemize}
    \item PyTorch 使用了更高效的数学库函数（如 \texttt{tanhf} 的向量化版本）
    \item PyTorch 的内核做了更细致的\textbf{向量化加载}——一次从内存读取多个连续元素
    \item PyTorch 内核可能使用了\textbf{更优的线程配置}和\textbf{寄存器分配}策略
\end{itemize}

\begin{table}[H]
\centering
\begin{tabular}{lcc}
\toprule
\textbf{实现} & \textbf{耗时 (ms)} & \textbf{说明} \\
\midrule
Manual（手写 PyTorch） & $\sim$8.1 & 多个内核，未融合 \\
CUDA（自写内核） & $\sim$1.8 & 单内核，已融合 \\
PyTorch（内置） & $\sim$1.1 & 单内核，高度优化 \\
\bottomrule
\end{tabular}
\caption{GeLU 三种实现的性能对比}
\end{table}

我们的 CUDA 实现已经比手写 PyTorch 快了 4.5 倍，但仍不如 PyTorch 内置版本。PyTorch 的实现做了更多底层优化（如向量化加载、更优的数学库函数）。

\begin{warningbox}{逐元素操作在 CUDA 中很简单，但...}
GeLU 是逐元素（elementwise）操作——每个输出元素只依赖一个输入元素，因此 CUDA 实现很直接。但大多数有趣的操作（矩阵乘法、softmax、RMSNorm）需要读取多个值，涉及\textbf{共享内存管理}、\textbf{线程同步}等复杂问题。
\end{warningbox}

\subsection{本章小结}

CUDA 让你直接控制 GPU 上的每个线程，实现内核融合。对于逐元素操作，CUDA 编程相对简单，主要是坐标计算和边界检查。但对于需要线程间通信的操作，CUDA 编程的复杂度会急剧上升——这正是 Triton 出现的原因。

%% ========== Section 7: Triton ==========

